\documentclass[11pt,a4paper]{article}
\usepackage[margin=2.2cm]{geometry}
\usepackage{fontspec}
\usepackage{booktabs}
\usepackage{array}
\usepackage{tabularx}
\usepackage{longtable}
\usepackage{hyperref}
\usepackage{xcolor}
\usepackage{enumitem}
\usepackage{titlesec}
\usepackage{fancyhdr}
\usepackage{graphicx}
\usepackage{amssymb}

\setmainfont{Times New Roman}
\definecolor{headingblack}{HTML}{000000}
\definecolor{lightgray}{HTML}{F2F4F7}
\hypersetup{colorlinks=true, linkcolor=headingblack, urlcolor=headingblack, citecolor=headingblack}
\titleformat{\section}{\Large\bfseries\color{headingblack}}{\thesection}{0.8em}{}
\titleformat{\subsection}{\large\bfseries\color{headingblack}}{\thesubsection}{0.8em}{}
\setlist[itemize]{leftmargin=1.4em}
\setlist[enumerate]{leftmargin=1.6em}
\pagestyle{fancy}
\setlength{\headheight}{14pt}
\fancyhf{}
\lhead{iBEC International Bioinformatics Engineering Competition}
\rhead{Project Report}
\cfoot{\thepage}

\begin{document}

\begin{center}
{\LARGE\bfseries OmicHub BioOps: Auditable Agent $\times$ Skill Infrastructure for Controlled Omics Analysis Delivery}\\[1.0em]
\end{center}

\begin{tabularx}{\textwidth}{>{\bfseries}p{0.24\textwidth}X}
\toprule
Team ID & \texttt{[TEAM-ID assigned by the organizing committee]}\\
Project Title & OmicHub BioOps: Auditable Agent $\times$ Skill Infrastructure for Controlled Omics Analysis Delivery\\
Team Name & \texttt{[Team name]}\\
Institution & OmicHub / Huazhong Agricultural University \texttt{[confirm official institution string]}\\
Team Members & Jian Zhang (Founder; product, architecture, full-stack development); \texttt{[add members and roles]}\\
Track / Category & \texttt{[If applicable]}\\
Date & 2026-09-18\\
\bottomrule
\end{tabularx}

\section*{Executive Summary}

OmicHub BioOps industrializes ``omics analysis delivery'', a high-frequency, long-chain task in life-science R\&D. The platform wraps expert-audited, version-controlled analysis pipelines into web-submittable tasks, and drives the whole process from requirement understanding to controlled execution with domain agents operating over a governed Skill library. Every analysis leaves a versioned pipeline definition, an environment snapshot, and a machine-verifiable evidence chain. Unlike general-purpose AI workbenches that let models generate workflows and trust their outputs, our agents only \emph{schedule} expert-audited pipelines: they hold no shell, pipeline-YAML, or database write permissions. Quality gates are rule-based (e.g., mapping rate $\geq 0.70$, Q30 $\geq 0.80$); missing rules force manual review. In the reference bulk RNA-seq scenario, a standard analysis submission drops from about 20 minutes of manual configuration and command-line operation to about 5 minutes through a validated web form ($-75\%$); incremental re-analysis within covered domains drops from about one expert-day to about 30 minutes ($-90\%$); pipeline customization drops from about five days to one day ($-80\%$). Project-level context documents, delivery READMEs, and MD5 verification lists turn ``unreproducible handover'' into ``reproducible takeover''. Platform v0.1.0 is in production use at the genomics laboratory of the College of Life Science and Technology, Huazhong Agricultural University.

\section{Project Background and Problem Definition}

Delivering one omics analysis in a pharma translational-medicine department, a breeding institute, a clinical-testing laboratory, or a university core facility typically spans data intake, sample/metadata validation, pipeline execution, quality control, statistical interpretation, and report delivery across several human roles. Today this collaboration is mediated by chat tools, spreadsheets, ad-hoc scripts, and task systems, which produces five structural problems:

\begin{itemize}
  \item \textbf{Fragmented entry points and error-prone submission.} One standard analysis requires manually preparing \texttt{config.yaml}, \texttt{samples.csv}, and \texttt{contrasts.csv} and submitting via the command line; errors surface only late in computation, wasting hours of compute.
  \item \textbf{Incremental requests depend on experts.} Follow-up questions on delivered results (a new contrast, an added enrichment analysis) must be re-parsed and re-designed by a bioinformatician; turnaround is measured in days.
  \item \textbf{High customization cost.} Swapping an aligner (STAR $\rightarrow$ HISAT2) or adding a QC step means rewriting and re-reviewing an entire pipeline, a multi-day to multi-week effort.
  \item \textbf{Untrustworthy AI artifacts.} Existing AI workbenches fully trust model-generated workflows and results; without reproducibility or traceability they cannot be accepted in compliance-heavy pharma and clinical-testing settings.
  \item \textbf{Organizational knowledge loss.} When analysts leave, their scripts and context are nearly impossible to reproduce; project archaeology from chat logs becomes the norm.
\end{itemize}

Problems 1--3 are efficiency issues, problem 4 is a trust issue, and problem 5 is an organizational-continuity issue. Existing tools each address one link; no mature commercial tool closes all three in one loop. The core engineering question is therefore: \emph{how can LLM-based agents exploit their reasoning flexibility without ever becoming the source of truth for execution or quality decisions?}

\section{Objectives and Scope}
\begin{enumerate}
  \item Primary objective: deliver a controlled omics-analysis platform in which domain agents decompose requirements, schedule \emph{expert-audited and versioned} pipelines, and produce machine-verifiable evidence chains (versioned pipeline YAML, environment snapshots, MD5-verified artifacts, audit events) for every run.
  \item Secondary objectives: (a) a governed Skill marketplace so incremental analysis needs are answered within declared capability boundaries without expert re-intervention; (b) module-level customization that inherits platform quality gates and delivery chains; (c) project-bound knowledge assets (AGENTS.md ledger, run-level README, artifact registry) that survive staff turnover.
  \item Scope boundary: the platform is a research/development analysis tool. It does not issue clinical diagnostic reports; GxP electronic-records compliance (21 CFR Part 11) adaptation is on the roadmap, not claimed as delivered. The multi-role collaboration layer is an experimental preview (disabled by default) and is not part of the claimed core path. Metrics marked as engineering estimates are pending pilot measurement.
\end{enumerate}

\section{Technical Route}

The end-to-end route is: form-validated data intake $\rightarrow$ preflight check $\rightarrow$ approval-gated task submission to a controlled queue $\rightarrow$ expert-audited pipeline execution (RNAFlow/Snakemake-based) $\rightarrow$ rule-based quality gates $\rightarrow$ evidence and artifact registration $\rightarrow$ delivery README with MD5 verification.

\begin{table}[h]
\centering
\renewcommand{\arraystretch}{1.25}
\begin{tabularx}{\textwidth}{p{0.18\textwidth}X X X}
\toprule
\textbf{Stage} & \textbf{Input} & \textbf{Method / Tool} & \textbf{Output}\\
\midrule
Data acquisition & FASTQ/BAM/sample sheets, project metadata & Web forms with parameter whitelists; read-only preflight Skill & Validated, project-bound input set\\
Analysis / modeling & Cleaned input + declared contrasts & Expert-audited pipelines (e.g., \texttt{rna\_seq} v2.3.0); 38 governed Skills; domain agents (LangGraph) & Counts, DEG/enrichment results, figures\\
Validation & QC metrics, artifacts & Rule-based gates (mapping $\geq 0.70$, Q30 $\geq 0.80$); MD5 \& artifact-integrity checks; manual-review fallback & Pass/fail/manual-review decisions + evidence events\\
Deployment / delivery & Validated results & Private on-premise deployment; delivery README + \texttt{environment.json} + \texttt{verification.tsv} & Reproducible delivery package\\
\bottomrule
\end{tabularx}
\caption{Technical route summary.}
\end{table}

\section{Methods and Implementation}

\begin{itemize}
  \item \textbf{Agent $\times$ Skill coupling.} Domain agents (e.g., RNA-seq analyst) understand requirements, decompose tasks, and explain results. Skills are stable business actions with explicit contracts: trigger condition + input schema + version + failure semantics + permission boundary (e.g., ``differential expression analysis v1.0.0''). Agents are assembled declaratively (identity / prompt / tool packs / capability boundary in \texttt{data/ai/*.yaml}), never by ad-hoc prompt stitching. The built-in Skill marketplace currently ships 38 executable Skills (36 analysis, 2 platform-management).
  \item \textbf{Expert-audited pipelines as versioned assets.} Built-in pipelines are constructed and reviewed by domain experts and stored as declarative, version-controlled YAML. Agents have no shell, no pipeline-YAML write access, and no database write access; submissions require an \texttt{approval\_token}, an \texttt{idempotency\_key}, and an \texttt{allowed\_flow}.
  \item \textbf{Evidence chain.} Key calls and state transitions emit evidence events bound to \texttt{trace\_id}, rule version, artifact references, and decision rationale. The controlled-execution base implements DAG runs, a plan-confirmation gate, an Artifact Registry, PostgreSQL Outbox + Redis Stream eventing with idempotent consumption, quality gates, and human-in-the-loop approval.
  \item \textbf{Knowledge and context.} A project-scoped storage root (\texttt{projects/\{slug\}/runs/\{task\}-\{timestamp\}/}) is enforced by platform code, not prompt convention. The AGENTS.md ledger and run-level delivery README (environment snapshot, dependency versions, file sizes, MD5 list) are auto-generated; a built-in verifier checks every delivered file (SUCCESS / MISSING / FAIL) and emits machine-readable \texttt{verification.tsv}. Knowledge-base RAG (PostgreSQL + pgvector, hybrid retrieval) serves domain SOPs with source- and version-attributed snippets.
  \item \textbf{Software environment.} Python $\geq$ 3.11 (FastAPI, SQLAlchemy 2 async, LangGraph), PostgreSQL (+pgvector), Redis, Vue 3 frontend, Snakemake-based RNAFlow pipelines in containerized execution; 29 application services and 227 backend Python modules under a strict DDD four-layer architecture; OpenTelemetry instrumentation for tracing.
  \item \textbf{Validation design.} Three end-to-end demonstration chains: (A) preflight-passed $\rightarrow$ approval $\rightarrow$ controlled submission $\rightarrow$ quality gate $\rightarrow$ delivery; (B) preflight-blocked (missing metadata stops before any compute is consumed); (C) quality-gate to manual review (below-threshold QC or missing rule version blocks delivery). Baseline comparison is the conventional manual workflow (command-line submission, expert-mediated re-analysis, full-pipeline rewrite).
\end{itemize}

\section{Engineering Workflow and Productization}

\begin{itemize}
  \item Requirement analysis: target users are pharma translational-medicine/bioinformatics teams, breeding and agricultural research groups, clinical-testing and research-service providers, and university core facilities; success criteria are submission-time, incremental-response time, and reproducibility of handover.
  \item Workflow design: four-layer architecture (workbench UI $\rightarrow$ Agent $\times$ Skill coupling layer $\rightarrow$ execution and knowledge base $\rightarrow$ infrastructure); a unified storage root shared by the tool pages, the AI assistant, and the agent workspace, so intermediate results flow by logical reference without physical copies.
  \item Implementation process: single-repository development with version control, DDD module boundaries (domain / application / API / infrastructure), declarative configuration for agents, Skills, and pipelines.
  \item Testing and iteration: unit and contract tests for the bridge/worker/MAS layers, idempotency and retry tests, preflight-block verification, and staging end-to-end acceptance runs; environment snapshots pin pipeline YAML versions and git commit SHAs.
  \item Deployment and maintenance: private on-premise deployment (data never leaves the customer network); containerized execution; platform v0.1.0 in production use at HZAU; roadmap includes a one-click deployment package and an Apache-2.0 open-source community edition.
\end{itemize}

\begin{table}[h]
\centering
\renewcommand{\arraystretch}{1.25}
\begin{tabularx}{\textwidth}{p{0.18\textwidth}X X X}
\toprule
\textbf{Process Stage} & \textbf{Key Activities} & \textbf{Evidence} & \textbf{Deliverable}\\
\midrule
Requirements & User/scenario analysis in HZAU genomics lab & Pain-point list, role/use-case matrix & Functional scope\\
Design & Four-layer architecture; Skill/pipeline contracts & Architecture diagram; module list & Technical design\\
Build & DDD implementation; 29 services, 227 modules; 38 Skills & Repository, commit history & Platform v0.1.0 prototype\\
Validate & Preflight / quality-gate / delivery chains & QC rule versions, verification.tsv, evidence events & Validated delivery package\\
Deliver & On-premise deployment, delivery README, MD5 list & Run archive + environment snapshot & Usable release\\
\bottomrule
\end{tabularx}
\caption{Engineering process summary.}
\end{table}

\section{Innovation Points}

\begin{itemize}
  \item \textbf{``AI schedules, never generates.''} Pipelines come from expert review and version control, not from model output; agents are confined to decomposition, scheduling, and explanation. This is the fundamental difference from generative-AI workbenches.
  \item \textbf{Machine-verifiable evidence chain.} Approval records, idempotent submissions, evidence events, environment snapshots, and MD5 verification together make every conclusion auditable by a checker, not by human trust.
  \item \textbf{Governed Skill marketplace.} 38 executable, contract-declared, version-managed Skills let agents answer incremental analysis needs within pre-declared capability boundaries (install-once, reuse-across-projects).
  \item \textbf{Code-enforced project memory.} Project-binding, ledgers, and delivery READMEs are generated by platform code rather than prompt convention, so organizational knowledge persists across staff turnover by construction.
  \item \textbf{Module-level customization.} Pipelines decompose into versioned modules (rule); a variant (STAR $\rightarrow$ HISAT2) inherits platform quality gates and delivery chains instead of re-validating the whole chain.
\end{itemize}

\section{Results and Evidence}

\begin{itemize}
  \item Result 1 (efficiency): in the reference bulk RNA-seq scenario, standard-analysis submission time drops from $\sim$20 min to $\sim$5 min ($-75\%$); expert-driven incremental re-analysis drops from $\sim$1 day to $\sim$30 min ($-90\%$, within covered Skills); pipeline customization drops from $\sim$5 days to $\sim$1 day ($-80\%$). Measurement basis: human operation and waiting time, excluding compute runtime; pilot-measured calibration planned.
  \item Result 2 (reproducibility): handover after staff turnover moves from ``not reproducible'' to ``reproducible'' via project context documents, run-level READMEs, MD5 checklists, and versioned pipeline YAML; a successor can reproduce an analysis and locate anomalies in about 10 minutes from the evidence chain.
  \item Result 3 (prototype in production): platform v0.1.0 serves the HZAU genomics laboratory, covering users, permissions, pipelines, tasks, files, and AI conversation (29 application services, 227 backend modules); preflight blocking prevents wasted computation; idempotent keys eliminate duplicate queueing; the Skill marketplace ships 38 executable Skills.
  \item Result 4 (controlled execution base): DAG orchestration, plan-confirmation gate, Artifact Registry, Outbox + Redis Stream eventing, quality gates, and HITL approval are implemented; the experimental multi-role collaboration layer is disabled by default and outside the core path.
\end{itemize}

\section{Application Value and Future Work}

Primary adopters are pharma translational-medicine departments, third-party clinical-testing and research-service providers, breeding/agriculture research teams, and university core facilities; the value proposition is standardized ``controlled task delivery'' with a complete evidence chain at private-deployment price points that replace outsourced-analysis budgets. The delivery model is horizontally transferable: preflight fields, allowed pipelines, artifact types, quality rules, and delivery templates are swappable configuration assets, so the same loop extends to ATAC-seq, single-cell, multi-omics, and eventually drug-screening or clinical-report-review tasks. Limitations: not all key parameters (e.g., aligner choice) are yet form-structured; quantified gains are engineering estimates pending pilot measurement; the structured SOP/QC rule vector store and data-flywheel auto-feedback are next-step enhancements. Next steps (12--18 months): three end-to-end demonstration chains at 2--3 pilot institutions, a one-click deployment package, the open-source community edition, and 2 software-copyright registrations plus 1 patent filing.

\section{Ethics, Safety, and Data Compliance}

\begin{itemize}
  \item Data privacy and licensing: private on-premise deployment; customer data physically never leaves the local network. Demo materials use de-identified or synthetic data only.
  \item Human genetic resources: services involving human genomic data comply with the PRC Regulations on the Administration of Human Genetic Resources and its 2023 Implementation Details (prior reporting, information backup, security review where applicable); the no-export architecture removes cross-border information-flow risk, and full audit logs/snapshots support customers' compliance filings.
  \item Clinical boundaries: the platform does not issue clinical diagnostic reports; it is scoped to research/development use, with customer obligations for clinical use written into commercial terms. GxP (21 CFR Part 11) adaptation is roadmap, with gap analysis scheduled for 2027; ISO 27001 program starts 2027.
  \item Responsible use: high-risk actions require human approval tokens scoped to one action/case/expiry; quality-rule gaps force manual review; agents cannot bypass gates or fabricate conclusions; encryption in transit and at rest.
\end{itemize}

\section{Published Outcomes}

\begin{itemize}
  \item Software release: OmicHub BioOps platform v0.1.0 (in production use at the genomics laboratory, College of Life Science and Technology, Huazhong Agricultural University); 38 built-in executable Skills and RNA-seq/ATAC-seq standard analysis pipelines.
  \item \texttt{[Add preprints, papers, patents, awards, or media reports if any; otherwise keep only the software release line.]}
\end{itemize}

\section{References}

\begin{enumerate}
  \item BCC Research, China genomics sequencing market data (USD 1.339B in 2020 $\rightarrow$ USD 4.235B in 2026E, CAGR $\approx$ 21.6\%), as cited in Novogene 2024 Annual Report.
  \item Novogene Co., Ltd. 2023 interim and 2024 annual reports (customer-base anchors).
  \item DNANexus/Seqera: Nextflow and Seqera Platform documentation.
  \item The Galaxy Project documentation, \url{https://galaxyproject.org}.
  \item DNAnexus, Seven Bridges, Terra platform documentation.
  \item LangGraph documentation, \url{https://langchain-ai.github.io/langgraph/}.
  \item E. Jones et al., ``pgvector: open-source vector similarity search for PostgreSQL,'' \url{https://github.com/pgvector/pgvector}.
  \item PRC State Council, Regulations on the Administration of Human Genetic Resources; MOST Implementation Details (effective 2023-07-01).
  \item U.S. FDA 21 CFR Part 11, Electronic Records; Electronic Signatures.
  \item Mendell et al.\ (Snakemake), ``Ten years of Snakemake,'' F1000Research 2024; RNAFlow internal pipeline suite.
\end{enumerate}

\section{Submission Checklist}
\begin{tabularx}{\textwidth}{p{0.08\textwidth}X X}
\toprule
\textbf{Check} & \textbf{Item} & \textbf{Notes}\\
\midrule
$\boxtimes$ & The report is mainly written in English. & English primary; Chinese notes omitted unless needed\\
$\square$ & Team ID is included on the title page and matches the submitted file names. & \texttt{[TEAM-ID]} placeholder --- fill and rename files\\
$\boxtimes$ & All figures and tables are numbered and explained. & Tables 1--2\\
$\boxtimes$ & Data sources, code availability, and software environment are documented. & Sections 4--5; see Supplementary Materials\\
$\boxtimes$ & Engineering workflow, testing evidence, and deployment or delivery status are described. & Sections 5--7\\
$\boxtimes$ & Published outcomes are separated from external references. & Sections 10--11\\
$\boxtimes$ & Claims match the technical abstract and supplementary materials. & Verified in drafting; re-check after edits\\
$\square$ & The final submitted file is exported as PDF. & Compile \texttt{.tex} with XeLaTeX, or export from Word\\
\bottomrule
\end{tabularx}

\end{document}
